% 子图 (c)：将通行 s 区间投至 ST 平面 ∩ 可达集 {t ≥ τ(s)}，得到走廊 Ω
% τ(s) = s/4 ⇒ ST 图上为过 (0,0) 与 (2.25, 9) 的直线
% 通行段：s ∈ [0, 6.5] ∪ [7, 9]；禁行切片：s ∈ [6.5, 7]
% Ω 下半段：(0,0) (1.625, 6.5) (3, 6.5) (3, 0)
% Ω 上半段：(1.75, 7) (2.25, 9) (3, 9) (3, 7)
% QP 轨迹严格起自 (0, 0)，紧贴 τ(s) 右下，止于禁行切片下沿 s ≈ 6.45
\begin{tikzpicture}
\begin{axis}[
    width=8.4cm, height=6.4cm,
    xlabel={$t$ (s)}, ylabel={$s$ (m)},
    xmin=0, xmax=3, ymin=0, ymax=9,
    grid=major, grid style={gray!18},
    axis line style={thick},
    tick label style={font=\footnotesize},
    label style={font=\small\bfseries},
    legend style={
        font=\scriptsize,
        legend columns=-1,
        column sep=0.9em,
        at={(0.5, -0.22)}, anchor=north,
        draw=none, fill=none,
    },
    legend cell align=left,
    legend image post style={line width=1.0pt},
]

% ===== 不可达区 t < τ(s)：(0,0) (2.25, 9) (0, 9) 三角 =====
\fill[gray!18]
    (axis cs:0, 0) -- (axis cs:2.25, 9) -- (axis cs:0, 9) -- cycle;

% ===== Ω 走廊 下半段 s ∈ [0, 6.5] =====
\fill[lightgreen!35]
    (axis cs:0, 0) -- (axis cs:1.625, 6.5)
    -- (axis cs:3, 6.5) -- (axis cs:3, 0) -- cycle;

% ===== Ω 走廊 上半段 s ∈ [7, 9] =====
\fill[lightgreen!35]
    (axis cs:1.75, 7) -- (axis cs:2.25, 9)
    -- (axis cs:3, 9) -- (axis cs:3, 7) -- cycle;

% ===== 禁行切片 s ∈ [6.5, 7] × 全 t =====
\fill[dangerred!25] (axis cs:0, 6.5) rectangle (axis cs:3, 7);
\draw[dangerred, thick, dashed] (axis cs:0, 6.5) -- (axis cs:3, 6.5);
\draw[dangerred, thick, dashed] (axis cs:0, 7  ) -- (axis cs:3, 7  );

% ===== τ(s) = s/4 蓝色虚线 =====
\addplot[deepblue, very thick, densely dashed] coordinates {(0, 0) (2.25, 9)};
\addlegendentry{$\tau(s) = s/4$}

% ===== Ω 走廊（图例占位用色块）=====
\addplot[lightgreen!35, fill=lightgreen!35, mark=none, forget plot] coordinates {(-1, -1)};
\addlegendimage{area legend, fill=lightgreen!35, draw=lightgreen!50!black}
\addlegendentry{$\Omega$ 走廊}

% ===== 禁行切片图例占位 =====
\addlegendimage{area legend, fill=dangerred!25, draw=dangerred}
\addlegendentry{禁行切片}

% ===== QP 最优轨迹 s(t) =====
\addplot[deeppink, very thick, -{Stealth}] coordinates {
    (0,    0)
    (0.15, 0.5)
    (0.32, 1.2)
    (0.55, 2.1)
    (0.80, 3.1)
    (1.05, 4.1)
    (1.30, 5.1)
    (1.55, 6.1)
    (1.70, 6.45)
};
\addlegendentry{QP 轨迹 $s(t)$}

% ===== 文字标注 =====
\node[lightgreen!35!black, font=\footnotesize\bfseries] at (axis cs:2.45, 3.0) {$\Omega$};
\node[dangerred!85!black, font=\tiny, anchor=west]
    at (axis cs:0.05, 8.55) {D 在 $t = \tau$ 时仍堵塞};
\node[dangerred, font=\scriptsize] at (axis cs:2.50, 6.75) {禁行切片};
\node[gray!55!black, font=\scriptsize, align=center]
    at (axis cs:0.55, 7.55) {$t < \tau(s)$\\不可达};

% ===== s 轴侧标 A、D 区间 =====
\draw[dangerred, thick] (axis cs:-0.04, 2.5) -- (axis cs:-0.04, 4);
\node[dangerred, font=\tiny, anchor=east] at (axis cs:-0.06, 3.25) {A};
\draw[dangerred, thick] (axis cs:-0.04, 6.5) -- (axis cs:-0.04, 7);
\node[dangerred, font=\tiny, anchor=east] at (axis cs:-0.06, 6.75) {D};

\end{axis}
\end{tikzpicture}
